% Generated draft: requires mathematical and source review.
\documentclass[11pt]{article}
\usepackage[margin=25mm]{geometry}
\usepackage{fontspec}
\setmainfont{Latin Modern Roman}
\usepackage{amsmath,amssymb,mathtools,unicode-math}
\setmathfont{Latin Modern Math}
\usepackage{enumitem,needspace,xurl,hyperref}
\setlength{\parindent}{0pt}
\setlength{\parskip}{5pt}
\setlength{\emergencystretch}{3em}
\newcommand{\N}{\mathbb{N}}
\newcommand{\Z}{\mathbb{Z}}
\newcommand{\Q}{\mathbb{Q}}
\newcommand{\R}{\mathbb{R}}
\newcommand{\C}{\mathbb{C}}
\newcommand{\F}{\mathbb{F}}
\newcommand{\A}{\mathbb{A}}
\newcommand{\PP}{\mathbb{P}}
\newcommand{\E}{\mathbb{E}}
\newcommand{\eps}{\varepsilon}
\newcommand{\dd}{\,\mathrm{d}}
\begin{document}
\section*{UnsolvedMath ID: 20001912}
\textbf{Source problem code:} catalogue entry
\subsection{Definitions and mathematical statement}
For a fixed closed smooth manifold \(M\) of dimension at least two, let \(\mathcal G^\infty(M)\) be its smooth Riemannian metrics. Interpret the source's entropy as topological entropy \(h_{\mathrm{top}}(\varphi_g)\) of the unit-speed geodesic flow, and put
\[
 \mathcal P=\{g\in\mathcal G^\infty(M):h_{\mathrm{top}}(\varphi_g)>0\}.
\]
Question \(8.4.1\) asks whether \(\mathcal P\) is dense in the \(C^k\) topology, for each specified \(k\), including \(k=\infty\). The excerpt does not itself define \(G\) or entropy; these choices are explicit interpretations. Dimension one must be excluded from the aid's all-dimensions assertion.

For \(S\subset\mathcal G^\infty(M)\), taking all closures within this smooth metric space, the aid states
\[
 \overline S^{\,C^\infty}=\bigcap_{r\geq0}\overline S^{\,C^r}.
\]
The smooth topology is generated by the finite derivative seminorms on the compact manifold. Consequently density in every finite \(C^r\) topology is equivalent to \(C^\infty\) density, but density in \(C^2\) alone implies neither \(C^3\) nor smooth density. The aid reports an explicit oscillatory metric-space example demonstrating this failure; its precise set is not reproduced in the supplied summary. Such a topology example makes no assertion that entropy behaves that way.

The proposed dynamical strategy is to create a nontrivial hyperbolic basic set, or horseshoe, by perturbations small in any prescribed finite \(C^r\) norm. Let \(\mathcal H\) be the smooth metrics whose geodesic flows contain such a positive-entropy hyperbolic basic set. Then \(\mathcal H\subset\mathcal P\); hyperbolic persistence makes this a robust certificate under sufficiently small \(C^2\) metric perturbations. In particular, if for every \(r\geq2\), every smooth \(g\), and every \(\eps>0\), one can find smooth \(\widetilde g\in\mathcal H\) with
\[
 \|\widetilde g-g\|_{C^r}<\eps,
\]
then \(\mathcal P\) is smoothly dense. Different finite orders may use different perturbations. The missing step is control of arbitrarily prescribed finite metric jets, not an abstract promotion of a \(C^2\) construction. This is a sufficient horseshoe-creation route; equivalence with all positive-entropy density mechanisms is not established here.

The source reports a positive \(C^2\) answer from [Contreras2010], also referring to [Con-Pat2002], and a positive \(C^\infty\) answer in dimension two from [Kni-Wei2002]. The aid qualifies the general \(C^2\) statement to dimensions at least two, describes the smooth surface cases as largely known, and reports finding no general higher-dimensional \(C^r\) theorem for \(r\geq3\). These are supplied status statements with incomplete theorem hypotheses and references, rather than a new literature determination or a proof of higher-topology density.
\subsection{Short English statement}
Positive-entropy density depends on the metric topology. Density in every finite derivative topology gives smooth density, but a \(C^2\) theorem alone does not. The aid proposes horseshoe creation with arbitrarily high finite derivative control as a sufficient route beyond the cited \(C^2\) and surface results.
\subsection{Sources}
\begin{itemize}
\item Local UnsolvedMath registry, permanent ID \(20001912\), source code \texttt{catalogue entry}; canonical collection \texttt{aim-geometry-notes.json}, notes[249]. Source item \(8.4\), actual question \(8.4.1\).
\item AIM workshop ``Geodesics,'' Burns--Matveev report ``Open problems and questions about geodesics'': \url{https://aimath.org/pastworkshops/geodesicsproblems.pdf}. The PDF was not opened in this drafting session.
\item Supplied reference labels: [Contreras2010] for \(C^2\) density, [Con-Pat2002] as a related reference, and [Kni-Wei2002] for the smooth surface case. Full bibliographic titles and exact theorem hypotheses are absent. The word ``recently'' belongs to the historical source.
\item Supplied research attempt \(1\), result ``reduction,'' run status \texttt{partially\_solved}, rigor label ``proved,'' novelty confidence ``low.'' Registry classification \texttt{PARTIAL-PROGRESS}, suggested difficulty \texttt{L4}, updated 14 August 2026. Original remarks and literature fields are empty; no proposer or year is supplied.
\end{itemize}
\end{document}
